\chapter{Scope and provenance}

\section{What this document is}
This is a byte-level reference for the System Exclusive implementation of the Technics
SX-WSA1R, reconstructed by reading the instrument's own firmware. Technics published a
user guide describing the \textsc{sysex bulk dump} feature and its error numbers, but no
MIDI implementation chart and no message format. Everything in the following chapters was
recovered from the program ROMs.

\section{Which firmware}
All addresses and all behaviour described here are those of the \textbf{OS v2.0} images:

\begin{center}
\begin{tabular}{llll}
\toprule
image & chip & factory part & role \\
\midrule
\texttt{prom\_a} & IC12 & QSIGCWSA1AX & CPU 1, MIDI and panel \\
\texttt{prom\_b} & IC13 & QSIGCWSA1BX & CPU 1, tables and display lists \\
\texttt{prom\_c} & IC28 & QSIGCWSA1CX & CPU 2 \\
\texttt{prom\_d} & IC21 & QSIGCWSA1DX & tone bank \\
\bottomrule
\end{tabular}
\end{center}

A v1 OS shipped and is not dumped. Nothing here should be assumed to hold for it.

\section{How to read the certainty markers}
Every substantive claim carries one of three markers.

\begin{description}
\item[\estab] Read directly from instructions or tables in the firmware. The symbol and
  address that establish it are cited.
\item[\likely] Follows from the code but was not executed or exhaustively checked.
\item[\unver] Stated in the firmware's own annotations or in Technics' documentation, but
  not confirmed here.
\end{description}

A wrong byte value in a document like this is worse than a gap, so where the evidence
stops, the text says so rather than rounding up to a plausible answer.

\section{One correction to earlier project notes}
\begin{tcolorbox}[colback=red!4!white,colframe=red!50!black,title=\textbf{The label
\sym{MIDI\_SysExHeader} does not do what its name suggests}]
\sym{MIDI\_SysExHeader} at \bytes{0xFA5CB8} contains the two bytes \bytes{F0 50}, and it is
tempting --- this project did it --- to cite it as the manufacturer-identifier test. It is
not. It is an \emph{outgoing} template, and no compare instruction anywhere in either
CPU~1 image reads that address. The acceptance decision is made by two immediate
comparisons inside \sym{MIDI\_RX\_SysExStart}; see section~\ref{sec:identifier}. The
resulting value, \bytes{0x50}, happens to be the same. The evidence is not.
\end{tcolorbox}
